\documentclass[10pt,a4paper]{amsart}
\usepackage[margin=19mm]{geometry}
\usepackage{amsmath,amssymb,mathtools}
\usepackage[scr=rsfs,cal=euler]{mathalfa}
\usepackage{xfrac}
\usepackage[unicode,hidelinks,bookmarks=false]{hyperref}
\newcommand{\ie}{i.e.\ }
\newcommand{\cf}{cf.\ }
\newcommand{\resp}{respectively\ }
\newcommand{\Subsection}{\subsection}
\providecommand{\nobreakdash}{\nobreak\hbox{-}}
\setlength{\emergencystretch}{2em}
\multlinegap0pt
\usepackage[normalem]{ulem}
\usepackage{amssymb}
\usepackage{xcolor}
\usepackage{float}
\usepackage{tikz}
\usepackage{booktabs}
\newtheorem{theorem}{Theorem}
\title{\textbf{A Discrete--Time Model of the Academic Pipeline in Mathematical Sciences with Constrained Hiring in the United States }}
\author{Oluwatosin Babasola, Olayemi Adeyemi, Ron Buckmire, Daozhou Gao, Maila Hallare, Olaniyi Iyiola, Deanna Needell, Chad M. Topaz, Andr\'es R. Vindas-Mel\'endez}
\date{Source: arXiv:2602.12188v1 (CC BY 4.0)}
\begin{document}
\maketitle

\noindent\emph{Source and licence.} \textbf{A Discrete--Time Model of the Academic Pipeline in Mathematical Sciences with Constrained Hiring in the United States }. Version arXiv:2602.12188v1 (CC BY 4.0), distributed by arXiv under the Creative Commons Attribution 4.0 International licence, http://creativecommons.org/licenses/by/4.0/, as recorded in the arXiv licence metadata for this exact version and shown on the abstract page https://arxiv.org/abs/2602.12188v1. Full licence notice: \texttt{licenses/arxiv-2602.12188-CC-BY-4.0.txt}.\par\smallskip

\noindent\textit{Editorial scope.} Counted unit: the article's theorem thm:boundedness (source0.tex lines 520--526). The article's complete original proof of it is delivered with it (source0.tex lines 528--558, one \texttt{proof} environment of 31 lines). No mathematics is written, repaired or re-derived: the statement and the proof are the article's own lines, in the article's own notation, and no retained line is substituted or re-flowed.  Retained uncounted as the source-local context the proof needs: lines 479--489.  Nothing was omitted from any retained span, so the package carries no omission record.  No retained line contains a citation, so the wrapper carries no bibliography and no bibliography entry.  No retained line contains a figure environment, an image-inclusion command or drawing code, so no image is delivered and the package declares no asset dependency.  Editorial formatting only, in addition: an AMS \texttt{amsart} preamble that restates the article's own declarations and macros used by the retained lines, namely the environments \texttt{theorem} on separate counters, together with the standard packages those lines need.  The retained environments print in this document as Theorem 1 in order of appearance, whereas the article prints the same items with the numbering of its own full document.  The complete native source of the article as distributed in the arXiv e-print for this exact version is bundled unchanged as \texttt{sources/194528/source0.tex}.
\begin{theorem}[Positivity and feasibility]\label{thm:positivity}
Assume \(U_0,G_0,P_0,F_0 \geq 0\) and \(B(t)\geq 0\) for all \(t\geq 0\). Suppose the coefficients satisfy
\[
g^U+a^U \leq 1,\qquad g^G+a^G \leq 1,\qquad p^{PF}(F_t)+a^P \leq 1 \quad \text{for all } t,
\]
and that the transition functions \(p^{UG}(G_t)\) and \(p^{PF}(F_t)\) are nonnegative for all \(t\).
Then the solution satisfies
\[
(U_t,G_t,P_t,F_t)\in\Omega \quad \text{for all } t\geq 0.
\]
\end{theorem}

\begin{theorem}[Boundedness under bounded inflow and positive exit]\label{thm:boundedness}
Assume \(0\leq B(t)\leq B_{\max}\) for all \(t\), and assume
\[
g^U+a^U>0,\qquad g^G+a^G>0,\qquad a^P>0,\qquad a^F>0.
\]
Under the constraints of Theorem~\ref{thm:positivity}, all trajectories \((U_t,G_t,P_t,F_t)\) are bounded for all \(t\geq 0\).
\end{theorem}

\begin{proof}
We derive comparison inequalities for each stock.

\noindent For undergraduates, we have
\[
U_{t+1} \leq (1-g^U-a^U)U_t + B_{\max}.
\]
Let \(c_U = 1-(g^U+a^U)\). Since \(g^U+a^U>0\), we have \(0\leq c_U<1\). Iteration yields
\[
U_t \leq U_0 + \frac{B_{\max}}{g^U+a^U},
\]
so \(U_t\) is bounded.

\noindent For graduates, we have
\[
G_{t+1} \leq (1-g^G-a^G)G_t + g^U U_t.
\]
Let \(c_G=1-(g^G+a^G)\) with \(0\leq c_G<1\). Since \(U_t\) is bounded, standard comparison implies boundedness of \(G_t\).

\noindent For postdocs, we have
\[
P_{t+1} \leq (1-a^P)P_t + p^{GP}g^G G_t.
\]
Since \(a^P>0\), let \(c_P=1-a^P\in[0,1)\). Boundedness of \(G_t\) implies boundedness of \(P_t\).

\noindent For faculty, we have
\[
F_{t+1} \leq (1-a^F)F_t + p^{GF}g^G G_t + P_t.
\]
With \(c_F=1-a^F\in[0,1)\) and bounded \(G_t,P_t\), it follows that \(F_t\) is bounded. Thus, all state variables are bounded.
\end{proof}

\end{document}
